\subsection{\textbf{Sprint 3}}

Ainda remoto, nossa Sprint 3 foi realizada durante toda a sua duração. Após a reunião com as \textit{stakeholders}, fizemos nossa \textit{planning} e foi decidido previamente que esta sprint, para o time de \textit{frontend}, seria basicamente para fazer a integração de todas as páginas e funcionalidades que tínhamos desenvolvido, com o \textit{backend}.

Minha \textit{squad} novamente ficou responsável por fazer a integração de módulos e cursos. Começamos pelos módulos, e tivemos um pouco de dificuldade no início para entender como deveriam ser feitas as requisições, que antes estavam \textit{mockadas}. Um problema que encontramos foram as mensagens de erros muito genéricas, o que não ajudou a identificar os problemas, exigindo testes extensivos para descobrir onde estavam os erros. Além disso, o Swagger apresentava muitas inconsistências entre o que era necessário ser passado e os exemplos fornecidos. Após várias tentativas, conseguimos fazer uma integração parcial de módulos e cursos.

Como o \textit{backend} não havia desenvolvido o \textit{endpoint} para que fosse possível o envio e requisição de imagens, isso ficou como débito técnico para a última sprint. Durante esta sprint, também criei o componente para utilizar na marcação de progresso do usuário conforme ele fosse concluindo os módulos.

Durante a apresentação, novamente apenas uma \textit{stakeholder} apareceu e, mais uma vez, não apontou nenhum ponto de melhoria ou algo que gostaria de ver diferente. Isso deixou o time frustrado, pois parecia que não estavam se importando com o nosso trabalho. Sinto também que, nesta sprint, não consegui me dedicar tanto ao projeto quanto gostaria.